% !PRINTED 612
\weberpage{612}

%<<CONT
according as the ideal $\mathfrak a$ is contained in the class $A_i$ or in some other
class. If we therefore multiply formula (2) by
$\chi_\varkappa(\overline{A_i'})$ and sum with respect to $\varkappa$, it follows that:

\weq{3}{
\frac{1}{h}\overset{\varkappa}{\sum}\chi_\varkappa(A_i') \log Q_k
% [sic] printed $\log Q_k$ with Latin $k$, but the summation index is $\varkappa$, so $\log Q_\varkappa$ is meant
= \overset{(1)}{\sum}\frac{1}{N(\mathfrak p)^s} + \frac{1}{2}\overset{(2)}{\sum}\frac{1}{N(\mathfrak p)^{2s}} + \frac{1}{3}\overset{(3)}{\sum}\frac{1}{N(\mathfrak p)^{3s}} + \cdots,
}
in which the sums $\overset{(1)}{\sum}, \overset{(2)}{\sum}, \overset{(3)}{\sum}, \ldots$ extend over all prime ideals $\mathfrak p$
whose first, second, third, $\ldots$ power is contained in the class $\overline{A_i}$.

We pass in (3) to the limit $s = 1$; the sums $\overset{(2)}{\sum}, \overset{(3)}{\sum}, \ldots$
retain a finite value in doing so, since their unconditional convergence does not cease for $s = 1$. The same holds of that part
of the sum $\overset{(1)}{\sum}$ which relates to prime ideals of higher than the
first degree, and so it remains questionable only whether
the sum

\[
\sum \frac{1}{N(\mathfrak p)^s},
\]

extended over all prime ideals of the first degree of the class $\overline{A_i}$, remains
finite or becomes infinite. The former would be the case if
there were in $\overline{A_i}$ no prime ideals of the first degree at all, or only
a finite number of them.

But if we assume the existence of a class field of degree $h$% faint pre-$h$ mark is scan/print damage, not a character
, then from § 164, 8. and § 163, (21) it follows that

\[
\log Q_2,\ \log Q_3,\ \ldots,\ \log Q_h
\]

retain finite limiting values, while $\log Q_1$ becomes infinite. The
left-hand side of (3) therefore becomes infinite, and the right-hand side cannot
remain finite.

Thus it is proved:

12. \gesperrt{If a class field exists whose degree
is not higher than the class number, then every
class contains infinitely many prime ideals of the first degree.}

\begin{center}
§ 166. \textbf{\gesperrt{Prime Ideals in the Ideal Classes.}}
\end{center}

If we take as group first of all the totality of the numbers $O$
of the field $\Omega$ without zero, then the classes are the ideal classes
 of the field $\Omega$, which, as we have seen, can be uniquely assigned to the classes
of quadratic forms of the discriminant $\varDelta$
%CONT>>
